\begin{afterword}
\summaryhead

The essay sets out to teach Bayes's theorem to readers who find it abstract. It opens with a screening problem: 1\% of women have breast cancer, a mammogram detects 80\% of cancers, and it gives false positives for 9.6\% of healthy women. Most doctors, the author says, overestimate the chance of cancer after a positive result. Counting women in a group of 10,000 gives the correct answer, 7.8\%.

A second problem, about eggs and pearls, repeats the method. The essay then introduces notation, the idea of degrees of freedom, the likelihood ratio, odds, and a scale of evidence in decibels taken from E. T. Jaynes. It derives the theorem from the worked examples.

The last third argues that Bayes's theorem underlies statistics, cognitive science and the scientific method, and that it explains and corrects Karl Popper's falsificationism. It ends by welcoming the reader into ``the Bayesian Conspiracy.''

\responsehead

The essay teaches one procedure: turn percentages into counts of people, then count. The procedure works, and it is Gigerenzer and Hoffrage's. Everything built on it is weaker.

Start with the evidence, since the essay's subject is how to weigh it. The opening tells the reader that ``only around 15\% of doctors'' solve the problem and that 46\% of doctors solve the natural-frequency version. Both figures match a cited study, Gigerenzer and Hoffrage (1995): 16\% and 46\%, for sixty students in Salzburg. A later study the essay does not cite found 46\% among physicians, but 10\%, not 15\%, with probabilities. The essay attributes to it a finding it does not contain, a ``slightly more evocative'' frequency format. When the reader asks whether the numbers are real, the essay answers that they have been ``extensively replicated'' and names no study.

The arithmetic fails where the lesson is sharpest. The first extreme case exists to show that none of the three numbers can be dropped, and it drops one: the answer is 1 in 125,000, not 1:100,000. The decibel method gives 83\% for a problem the essay solved as 84\% a page earlier. The assumptions that make the method valid are stated wrongly or not at all. Tests are called ``independent'' when they must be independent given the disease. Odds are called an ``odds ratio.''

The essay announces that it has ``entirely succeeded'' at 68\% of its length and then changes its business. It stops teaching and starts recruiting. A ``Bayesian revolution'' is said to be sweeping the sciences, and no scientist is named. Popper is said to hold that theories can be ``definitely falsified.'' He held that a refutation is decisive in logic, and wrote that in practice ``no conclusive disproof of a theory can ever be produced.'' Newton's theory is said to have fallen to a hidden fact. In fact the anomaly in Mercury's orbit was known for more than fifty years before Einstein explained it, and Newtonians kept their theory throughout. The theorem is said to put ``physical causality'' on its right side, which the essay's own eggs refute.

The one question on which the claim about science depends, where the prior for a theory comes from, is answered with jokes about Kazaa and the ``Bayes Council.'' The serious answer that follows covers only base rates that can be counted. The frame runs from ``Soon you will be one of us'' to ``You are now an initiate.''

In short: a borrowed counting method, wrapped in misreported studies, a miscalculated example and a misdescribed Popper, used to enlist the reader in a revolution the essay never shows exists.
\end{afterword}
